% T008: reference solution
循环在 $P$-$V$ 平面顺时针，气体对外净做功等于矩形面积：
\[
W_{\rm net}=(2P_0-P_0)(2V_0-V_0)=P_0V_0.
\]
对单原子理想气体，$C_V=\tfrac32R$、$C_P=\tfrac52R$。由 $PV=nRT$，四点温度依次为
\[
T_A=\frac{P_0V_0}{nR},\quad
T_B=2T_A,\quad T_C=4T_A,\quad T_D=2T_A.
\]
吸热只发生在 $A\to B$ 的定容升温与 $B\to C$ 的定压膨胀阶段：
\[
Q_{AB}=nC_V(T_B-T_A)=\frac32P_0V_0,
\qquad
Q_{BC}=nC_P(T_C-T_B)=5P_0V_0.
\]
因此 $Q_{\rm in}=13P_0V_0/2$，循环热效率
\[
\eta=\frac{W_{\rm net}}{Q_{\rm in}}=\frac{2}{13}.
\]
其余两段放热，净热量等于净功，符合循环内能变化为零的第一定律。
